% !TEX root = ../article.tex
% =============================================================================
% Macros propias y metadatos del PDF
% =============================================================================

% --- Metadatos del PDF ---
\hypersetup{
    pdftitle={\articleTitle},
    pdfauthor={\articleAuthorsPDF},
    pdfsubject={\articleDocument{} - \articleProgram, \articleInstitution},
    pdfkeywords={\articleKeywords},
    pdflang={es-PE}
}

% --- Fuente de figuras y tablas ---
% Debajo de cada figura o tabla:
%   \articleSource{Elaboración propia.}
%   \articleSource{Adaptado de \textcite{clave}.}
\newcommand{\articleSource}[1]{%
    \par\vspace{0.3em}%
    {\footnotesize\centering\textit{Fuente:} #1\par}%
}

% --- Notas mientras se escribe ---
% \pendiente{...} marca en rojo lo que falta. Antes de entregar no debe
% quedar ninguno: los scripts de compilación avisan cuántos hay.
\newcommand{\pendiente}[1]{\textcolor{red}{[PENDIENTE: #1]}}

% --- Resumen ejecutivo de una página ---
% La guía pide que no exceda una página. sections/01-executive-summary.tex
% pone \label{summary:start} y \label{summary:end}; si caen en páginas
% distintas, LaTeX emite un aviso que los scripts de compilación muestran.
% El aviso dice "hoja" y no "página" a propósito: pdflatex no escribe bien
% las tildes en el .log, y los scripts buscan este texto exacto.
\newcommand{\articleCheckOnePage}{%
    \ifnum\getpagerefnumber{summary:end}>\getpagerefnumber{summary:start}\relax
        \PackageWarningNoLine{thoraxis}{El resumen ejecutivo excede una hoja}%
    \fi
}
